\SourceLeaf{361}%
nung vorgenommen werden soll, sehr genäherte Werthe der zu bestimmenden Grössen schon vorliegen. Genau genommen haben übrigens auch alle andern Auflösungsarten des Problems, namentlich auch die \textsc{Bessels}che, theilweise diesen Charakter indirecter Operationen. Der wesentlichste Umstand bleibt aber der, dass von den wiederholten Annäherungen nur die letzte, die den ganzen Kern der Rechnung vollständig enthält, aufbewahrt zu werden braucht, und dass diese eine Kürze und Übersichtlichkeit hat, wie keine andere Methode.

Die Formeln für die Auflösung der Aufgabe auf der Sphäroidfläche unterscheiden sich von denen für die Kugelfläche lediglich dadurch, dass gewisse Zwischengrössen, die bei diesen constant sind. bei jenen von der Breite abhängig werden; diese lassen sich folglich in eine Hülfstafel bringen, deren Argument die Breite bildet. Steht eine solche Hülfstafel zu Gebote, so wird in jedem concreten Falle die Rechnung auf der Sphäroidfläche ganz eben so leicht, wie auf der Kugelfläche. Für die Zone von 51---54 Grad Breite, welche für das Hannoversche Dreieckssystem ausreicht, ist eine solche Hülfstafel am Schlusse der Abhandlung beigefügt, und zwar nach demjenigen Werthe der Abplattung, welchen \textsc{Bessel} aus allen bisherigen Gradmessungen abgeleitet hat, und der auch schon in der ersten Abhandlung zum Grunde gelegt war. Wer dieselbe Methode auf ein ausserhalb dieser Zone liegendes Dreieckssystem anwenden wollte, würde damit anfangen müssen, jene Hülfstafel für seinen Zweck weiter auszudehnen, oder, falls er eine andere Abplattung zum Grunde legen wollte, sich erst eine neue Hülfstafel zu construiren. Wo es die Bearbeitung eines grossen Dreieckssystems gilt, kommt eine solche vorgängige Hülfsarbeit gar nicht in Betracht, und die darauf gewandte Mühe wird durch die Bequemlichkeit der Benutzung reichlich ersetzt. Für den Fall hingegen, wo man nur eine oder ein paar concrete Auflösungen der in Rede stehenden Aufgabe suchen soll, hat die Methode nicht vorzugsweise bestimmt sein sollen.
